\qquad A necessidade de trabalhar com dados e informações ordenados está presente em diversas situações do cotidiano. Ao procurar um contato em uma lista telefônica, por exemplo, a busca seria mais complexa caso os nomes não estivessem organizados em ordem alfabética. O mesmo vale para a listagem de produtos por preço em uma plataforma de e-commerce ou para a classificação de alunos por nota em um sistema acadêmico.

\qquad Diante disso, fica claro que as operações de ordenação desempenham um papel central na Ciência da Computação. A ordenação de dados é uma tarefa fundamental que impacta diretamente a eficiência de algoritmos de busca, a otimização de consultas em bancos de dados e a performance no processamento de grandes volumes de informação.

\qquad O objetivo principal deste relatório é analisar e comparar o comportamento e o tempo de execução de cinco algoritmos de ordenação: \textit{Bubble Sort}, \textit{Insertion Sort}, \textit{Selection Sort}, \textit{Shell Sort} e \textit{Merge Sort}. Os três primeiros possuem complexidade quadrática e são utilizados com fins didáticos. O \textit{Shell Sort} apresenta desempenho subquadrático, e o \textit{Merge Sort}, único algoritmo recursivo do trabalho, baseia-se na estratégia de dividir para conquistar e possui complexidade $O(n \log n)$ em todos os casos. Esse conjunto permite contrastar diferentes classes de eficiência e diferentes estratégias de ordenação.

\qquad A análise de cada algoritmo é conduzida sob três cenários de entrada: vetores ordenados em ordem crescente (melhor caso), em ordem decrescente (pior caso) e com dados aleatórios (caso médio). Para cada cenário foram utilizados vetores de 10, 100, 1.000, 10.000, 100.000 e 1.000.000 de elementos, totalizando dezoito execuções por algoritmo. As implementações foram desenvolvidas em C++ e os tempos de execução registrados em segundos.

\qquad O trabalho está organizado da seguinte forma: primeiro é apresentada a implementação de cada algoritmo, acompanhada da tabela de custo de suas instruções; em seguida, a análise de complexidade com teste de mesa e desenvolvimento matemático para os três cenários, incluindo a resolução da recorrência no caso do \textit{Merge Sort}; e, por fim, os resultados experimentais, com tabelas e gráficos individuais de cada algoritmo e uma seção comparativa reunindo os cinco. Essa estrutura permite confrontar os resultados práticos obtidos experimentalmente com a teoria da complexidade de algoritmos, identificando as vantagens, desvantagens e cenários de uso adequados para cada método estudado.